% Results

\section{Results}

We demonstrate our party-specific overlapping district system through two test cases: Ohio (15 congressional seats) and California (52 seats). These states represent moderate-sized and large delegations, respectively, with different partisan compositions.

\subsection{Ohio: Moderate-Sized Delegation}

Ohio's 15 congressional seats provide an ideal test case for our system. Using 2020 presidential election results (Biden 45.3\%, Trump 53.2\%, Other 1.5\%), we apply the D'Hondt method with the natural threshold of $1/(N+1) = 6.25\%$.

\paragraph{Seat Allocation}
The proportional allocation yields 7 Democratic seats and 8 Republican seats, closely matching the statewide vote shares (46.7\% vs 53.3\% of seats for actual votes of 45.3\% vs 53.2\%). The ``Other'' category falls below the threshold and receives no seats.

\paragraph{Party-Weighted Districts}
We generate two independent sets of overlapping districts---one for Democrats, one for Republicans---each covering the entire state. Critically, districts are balanced by party-vote weighted population rather than total population:

\begin{itemize}
    \item \textbf{Democratic districts}: 7 districts, each containing approximately 763,000 Democratic-weighted residents (Ohio's 11.8M population $\times$ 45.3\% $\div$ 7 districts)
    \item \textbf{Republican districts}: 8 districts, each containing approximately 785,000 Republican-weighted residents (11.8M $\times$ 53.2\% $\div$ 8 districts)
\end{itemize}

This weighting ensures equal representation \emph{within} each party's caucus. Every Democratic voter has an equal say in selecting Democratic representatives, regardless of whether they live in Cleveland or rural Appalachia.

\paragraph{Comparison to Traditional Approach}
Under traditional single-member districts balanced by total population, each district would contain approximately 787,000 residents. However, this creates severe imbalances in party representation. A Democratic district in urban Cleveland might contain 560,000 Democratic voters, while a Democratic district forced into rural areas might contain only 330,000---a 1.7$\times$ imbalance that dilutes rural Democratic voters' representation within the Democratic caucus.

Our party-weighted approach eliminates this imbalance entirely. Table~\ref{tab:ohio-balance} shows the variation in Democratic votes across the 7 Democratic districts ranges only 392,725 to 393,386 (0.17\% variation)---effectively perfect balance.

\begin{table}[h]
\centering
\caption{Ohio Democratic districts: Party-vote balance}
\label{tab:ohio-balance}
\begin{tabular}{lrrr}
\toprule
District & Total Pop. & Dem. Votes & Deviation \\
\midrule
District 1 & 1,684,925 & 392,928 & -0.04\% \\
District 2 & 1,685,847 & 393,140 & +0.01\% \\
District 3 & 1,686,321 & 393,259 & +0.04\% \\
District 4 & 1,686,871 & 393,386 & +0.07\% \\
District 5 & 1,685,594 & 393,098 & 0.00\% \\
District 6 & 1,685,864 & 393,147 & +0.01\% \\
District 7 & 1,684,026 & 392,725 & -0.10\% \\
\midrule
\textbf{Average} & 1,685,635 & 393,098 & --- \\
\textbf{Std. Dev.} & 1,012 & 226 & 0.06\% \\
\bottomrule
\end{tabular}
\end{table}

\subsection{California: Large Delegation with Third-Party Representation}

California's 52-seat delegation demonstrates the system's scalability and its capacity to represent minor parties. With 2020 results (Biden 63.4\%, Trump 34.3\%, Other 2.3\%), the natural threshold of $1/(52+1) = 1.89\%$ allows the ``Other'' category to qualify for representation.

\paragraph{Seat Allocation}
The D'Hondt allocation yields:
\begin{itemize}
    \item 33 Democratic seats (63.5\% of seats for 63.4\% of votes)
    \item 18 Republican seats (34.6\% of seats for 34.3\% of votes)
    \item 1 Other seat (1.9\% of seats for 2.3\% of votes)
\end{itemize}

\paragraph{Three-Party Overlapping Districts}
California now has \emph{three} independent district maps covering the state:

\begin{itemize}
    \item \textbf{33 Democratic districts}: Average 759,613 Democratic-weighted residents per district
    \item \textbf{18 Republican districts}: Average 753,422 Republican-weighted residents per district
    \item \textbf{1 Other district}: Represents all 39.5M California residents for third-party voters (at-large)
\end{itemize}

The single ``Other'' seat functions as an at-large district representing the entire state, giving voice to Libertarians, Greens, and independent voters who would otherwise have no representation.

\paragraph{Proportionality Achievement}
Table~\ref{tab:proportionality} compares our system to California's actual 2022 congressional delegation. Our system achieves near-perfect proportionality (efficiency gap $< 0.1\%$) by construction, while the actual districts exhibit typical partisan bias.

\begin{table}[h]
\centering
\caption{Proportionality comparison: Our system vs. actual California districts}
\label{tab:proportionality}
\begin{tabular}{lcccc}
\toprule
System & Vote \% & Seats & Seat \% & Eff. Gap \\
\midrule
\textbf{Our System (2020)} \\
\quad Democratic & 63.4\% & 33 & 63.5\% & \multirow{3}{*}{0.08\%} \\
\quad Republican & 34.3\% & 18 & 34.6\% & \\
\quad Other & 2.3\% & 1 & 1.9\% & \\
\midrule
\textbf{Actual Districts (2022)} \\
\quad Democratic & --- & 42 & 80.8\% & \multirow{2}{*}{15-20\%} \\
\quad Republican & --- & 10 & 19.2\% & \\
\bottomrule
\end{tabular}
\end{table}

\subsection{Key Findings}

Our results demonstrate three critical properties of party-specific overlapping districts:

\paragraph{Perfect Internal Balance}
Within each party's delegation, every representative serves an equal number of their party's voters (standard deviation $< 0.2\%$). This contrasts sharply with traditional districts, where geographic constraints force party representatives to serve vastly different numbers of co-partisans depending on urban/rural geography.

\paragraph{Guaranteed Proportionality}
By allocating seats \emph{before} drawing districts, our system achieves proportionality by construction. The efficiency gap approaches zero, limited only by rounding constraints in seat allocation.

\paragraph{Third-Party Representation}
In large delegations, the natural threshold becomes low enough for significant third parties to win representation. California's 1.89\% threshold enabled ``Other'' voters to elect a representative despite comprising only 2.3\% of the electorate---impossible under traditional single-member plurality rules.
